% =====================================================================
%  Acquired capabilities and the extension path
%  SCOPE: Phase A/B subsystems not exercised by Phase C, and what each
%         would cost to bring back into the delivered system.
% =====================================================================
\section{Acquired Capabilities and the Extension Path}
\label{sec:extension}

\subsection{What this section is}

\phaseC{} delivers a connected measurement node. Doing so exercises a
strict subset of what was built and measured in Phases~A and~B: the
platform, the twin, the safety layer, the lidar, the cameras, the
regression discipline. Several substantial subsystems --- mapping, SLAM,
graph search, fruit detection, the five-DOF arm --- are \emph{present in
the repository, validated, and not invoked by the delivered mission}.

That is a design outcome, not an omission, and it deserves stating
precisely because the two readings are easy to confuse:

\begin{itemize}[leftmargin=*,itemsep=1pt,topsep=2pt]
  \item These subsystems were \textbf{not} planned and left unbuilt.
        They were built, instrumented and measured, and their results
        --- including two negative ones --- are reported in
        Sections~\ref{sec:phaseA_results} and~\ref{sec:results}.
  \item They are \textbf{not} dormant code of unknown quality. Each is
        covered by the \phaseB{} suite, and the \phaseC{} suite asserts
        that \phaseC{} never writes into the \phaseB{} tree.
  \item They are \textbf{not} required by the delivered mission, because
        the mission operates in a catalogued space
        (Sec.~\ref{sec:phaseC_catalogue}).
\end{itemize}

They constitute an \emph{extension path}: the capabilities the platform
already has, which a wider deployment would need, together with what
each would cost to bring back in.

\subsection{The capabilities, and what each buys}

\begin{table*}[!t]
\caption{Capabilities built and measured in Phases~A--B, their status
against the \phaseC{} mission, and the condition under which each stops
being optional. ``Measured'' refers to the results in
Sec.~\ref{sec:results}.}
\label{tab:extension}
\centering\footnotesize
\renewcommand{\arraystretch}{1.25}
\begin{tabular}{@{}p{0.17\textwidth}p{0.24\textwidth}p{0.24\textwidth}p{0.25\textwidth}@{}}
\toprule
\textbf{Capability} & \textbf{Status after Phases A--B} &
\textbf{Why \phaseC{} does not use it} & \textbf{What would require it} \\
\midrule

Occupancy mapping \newline (log-odds, 2D lidar) &
\ok{Works.} Interior reconstructed to
\measured{$-4$~cm~/~$+10$~cm}, \measured{100\%} floor coverage,
\measured{0.0\%} phantom clutter. &
The greenhouse is generated from a catalogue that predates the robot; a
map would be a lossier copy of a file already on disk. &
A greenhouse whose layout is \emph{not} surveyed, or one where gutters
are moved between seasons without the catalogue being updated. \\

A$^\star$ path planning &
\ok{Works}, and was tuned through inflation and clearance problems in a
\SI{0.16}{m} gap. &
Free space is four corridors joined at both ends; the closed-form route
cannot return a novel answer, which in that clearance is a feature
(Sec.~\ref{sec:phaseC_nav}). &
Mobile obstacles --- a trolley, a person, a second robot --- or any
layout not decomposable into parallel bands. \\

Correlative scan-matching SLAM &
\bad{Measured worse} than the calibrated odometry prior it was meant to
correct; disabled in the production configuration
(Sec.~\ref{sec:slam}). &
Not used, and would not be used even if it worked: the visual mark
alignment closes position against a physical feature of the building,
which is stronger than either. &
Nothing, as implemented. A loop-closing back-end would be a new
subsystem, not a re-enabling of this one. \\

Colour fruit detection \newline and map placement &
\warn{Partially works.} \measured{77\%} recall, \measured{0.22~m}
localisation error, roughly two of three map estimates spurious. &
\phaseC{} photographs the plant for a human and does not classify it.
The detector's error is a quarter of the plant spacing --- enough to
identify the plant, not the berry. &
Any autonomous decision taken \emph{from} the image: ripeness triage,
disease flagging, yield estimation. \\

Five-DOF arm and gripper &
\ok{Works.} Pick-and-place sequence validated in both simulators. &
The mission measures and photographs; it does not touch the plant. &
Sampling that requires contact --- leaf collection, substrate probing,
selective harvesting. \\

Behaviour-tree mission \newline supervision &
\ok{Works.} Multi-trip harvesting with dynamic re-planning. &
A \phaseC{} mission is a list of stations with a defined end. A queue
and a loop express that exactly; a behaviour tree would express it
loosely. &
Missions with conditional branches: recharge when low, re-measure on an
anomaly, defer a blocked station and return to it. \\

\bottomrule
\end{tabular}
\end{table*}

Table~\ref{tab:extension} is also the answer to the question a reader
may reasonably ask on finishing Sec.~\ref{sec:phaseC_nav}: the platform
did not lose the ability to navigate an unknown space. It retains it,
and the delivered mission declines to use it because the space is known.

\subsection{The extension that the measured results argue against}

One row of Table~\ref{tab:extension} deserves emphasis, because it is
the one place where ``we already built it'' would be the wrong
conclusion to draw.

The homemade SLAM front-end was implemented, instrumented, and found to
be \emph{worse} than the dead-reckoning prior it corrects
(Sec.~\ref{sec:slam}). \phaseC{} then solved the same underlying problem
--- knowing precisely where the robot is at the moment it takes a reading
--- by a different route entirely: it does not improve the global pose
estimate at all. It accepts that odometry drifts, and closes the last
centimetres visually against a painted physical feature whose position
is known exactly (Sec.~\ref{sec:phaseC_align}).

That comparison is worth more than either result alone. A global
localisation solution accurate to \measured{\SI{0.20}{m}} does not
support a task whose labels are \SI{0.10}{m} apart, and no incremental
improvement to scan matching was going to change that by the required
order of magnitude. A local, absolute, fiducial-based correction does
support it --- and it is cheaper, simpler, and verifiable from a single
image. Where a task can be re-posed so that global accuracy stops
mattering, that is usually the better engineering.

\subsection{Priority, if the work continues}

Ordered by value against cost for this deployment, and drawing on the
finding of Sec.~\ref{sec:phaseC_route} that per-station dwell time now
dominates the energy budget:

\begin{enumerate}[leftmargin=*,itemsep=2pt,topsep=3pt]
  \item \textbf{Replace the synthesised field with real probes.} This is
        the largest gap between the delivered system and a deployable
        one, and it is one function
        (Sec.~\ref{sec:phaseC_measurement}). Everything downstream ---
        encoding, sealing, transport, storage, display --- is already
        exercised end to end.
  \item \textbf{Measure $P_{\text{idle}}$ and $P_{\text{drive}}$ on the
        bench.} Two numbers turn Eq.~\ref{eq:energy} from a stated model
        into a measured budget, and the autonomy figure a fleet operator
        actually needs follows immediately.
  \item \textbf{Attack per-station dwell time}, which is where the
        remaining energy is: the visual settling tolerance, the number of
        correction attempts, and image encoding cost.
  \item \textbf{Re-enable obstacle handling} (A$^\star$ over a local
        costmap seeded from the catalogue) before any deployment sharing
        aisles with people. The brake already stops for an unexpected
        object; it has nowhere to route around one.
  \item \textbf{Fixed ESP nodes.} The protocol already carries
        \param{node\_kind} and namespaces every topic per node
        (Sec.~\ref{sec:phaseC_protocol}); the Cloud already keeps nodes
        in a dictionary. This is the cheapest capability increase
        available, because the interface for it was built and tested
        before there was anything to plug into it.
\end{enumerate}
